\documentclass[11pt,a4paper]{article}
\newcommand{\fichetype}{Corrigé détaillé}
% Thème A4 commun, librement modifiable.
\usepackage[T1]{fontenc}
\usepackage[utf8]{inputenc}
\usepackage[french]{babel}
\usepackage{lmodern,amsmath,amssymb,mathtools,tikz,pgfplots}
\usepackage[margin=2.2cm,headheight=15pt]{geometry}
\usepackage{xcolor,fancyhdr,enumitem,hyperref,needspace}
\usetikzlibrary{arrows.meta,positioning,calc}
\pgfplotsset{compat=1.18,width=11cm,height=6cm}
\definecolor{Ink}{HTML}{202124}
\definecolor{Blue}{HTML}{356B9E}
\definecolor{Muted}{HTML}{737579}
\color{Ink}
\renewcommand{\familydefault}{\sfdefault}
\setlength{\parindent}{0pt}
\setlength{\parskip}{6pt}
\setlist{itemsep=3pt,topsep=4pt}
\pagestyle{fancy}\fancyhf{}
\fancyhead[L]{\small\color{Muted}MATHS LYCÉE}
\fancyhead[R]{\small\color{Muted}2026--2027}
\fancyfoot[L]{\small\color{Muted}\fichetype}
\fancyfoot[R]{\small\thepage}
\renewcommand{\headrulewidth}{0.3pt}
\hypersetup{colorlinks=true,urlcolor=Blue,linkcolor=Blue}
\newcommand{\R}{\mathbb{R}}
\newcommand{\N}{\mathbb{N}}
\newcommand{\E}{\mathbb{E}}
\newcommand{\V}{\mathrm{V}}
\newcommand{\exercice}[3]{\Needspace{6\baselineskip}\vspace{0.5cm}{\color{Blue}\small\textbf{EXERCICE #1}\hfill #2}\par{\large\bfseries #3}\par}

\begin{document}
{\small\color{Blue}PREMIÈRE SPÉCIALITÉ}\par{\LARGE\bfseries Second degré : formes, équations et signes}\par{\large Corrigé détaillé}\par\vspace{0.2cm}\hrule\vspace{0.4cm}
\par\noindent\begin{minipage}{\linewidth}\exercice{1}{Application}{Reconnaître et convertir}
{\small\color{Muted}\textbf{Énoncé.} Pour $f(x)=x^2-8x+7$, déterminer les coefficients, compléter le carré et donner le sommet. En déduire le minimum.}
\par\medskip
\textbf{Résolution.}\par
Les coefficients sont $a=1$, $b=-8$, $c=7$. Comme $(x-4)^2=x^2-8x+16$, on a $f(x)=(x-4)^2-9$. Le sommet est $(4,-9)$. Le carré étant positif ou nul, $f(x)\ge-9$, avec égalité pour $x=4$ : le minimum est $-9$.
\end{minipage}\par
\par\noindent\begin{minipage}{\linewidth}\exercice{2}{Application}{Trois discriminants}
{\small\color{Muted}\textbf{Énoncé.} Résoudre dans $\mathbb R$ : $x^2-7x+10=0$ ; $4x^2+4x+1=0$ ; $x^2+2x+5=0$.}
\par\medskip
\textbf{Résolution.}\par
Pour la première, $\Delta=49-40=9$, d’où $x=(7\pm3)/2$ : $2$ et $5$. Pour la deuxième, $\Delta=16-16=0$, donc $x=-4/8=-1/2$ ; on retrouve $(2x+1)^2=0$. Pour la troisième, $\Delta=4-20=-16<0$ : aucune solution réelle.
\end{minipage}\par
\par\noindent\begin{minipage}{\linewidth}\exercice{3}{Entraînement}{Choisir la factorisation}
{\small\color{Muted}\textbf{Énoncé.} Factoriser $9x^2-16$, $x^2+x-12$ et $3x^2-12x$. Donner leurs racines.}
\par\medskip
\textbf{Résolution.}\par
Une différence de carrés donne $9x^2-16=(3x-4)(3x+4)$, racines $\pm4/3$. Les nombres $3$ et $-4$ ont pour somme $-1$ et produit $-12$, donc $x^2+x-12=(x-3)(x+4)$, racines $3$ et $-4$. Enfin $3x^2-12x=3x(x-4)$, racines $0$ et $4$.
\end{minipage}\par
\par\noindent\begin{minipage}{\linewidth}\exercice{4}{Entraînement}{Signe et inéquations}
{\small\color{Muted}\textbf{Énoncé.} Résoudre $2x^2-5x-3\le0$, puis $-x^2+2x+3<0$.}
\par\medskip
\textbf{Résolution.}\par
Le premier discriminant vaut $25+24=49$. Les racines sont $(5-7)/4=-1/2$ et $(5+7)/4=3$. Comme $a>0$, le trinôme est négatif entre elles : solution $[-1/2,3]$. Pour le second, $-x^2+2x+3=-(x-3)(x+1)$. Son signe est négatif à l’extérieur des racines. L’inégalité est stricte : $]-\infty,-1[\cup]3,+\infty[$.
\end{minipage}\par
\par\noindent\begin{minipage}{\linewidth}\exercice{5}{Entraînement}{Une parabole imposée}
{\small\color{Muted}\textbf{Énoncé.} Un trinôme s’annule en $2$ et $5$ et vaut $-20$ en $0$. Déterminer son expression. Calculer son extremum.}
\par\medskip
\textbf{Résolution.}\par
Il s’écrit $f(x)=a(x-2)(x-5)$. En $0$, $10a=-20$, donc $a=-2$. Ainsi $f(x)=-2x^2+14x-20=-2(x-7/2)^2+9/2$. Le carré étant positif, $f(x)\le9/2$. Le maximum vaut $9/2$, atteint en $7/2$.
\end{minipage}\par
\par\noindent\begin{minipage}{\linewidth}\exercice{6}{Entraînement}{Clôturer un jardin}
{\small\color{Muted}\textbf{Énoncé.} On dispose de $24$ m de grillage pour trois côtés d’un rectangle, le quatrième bordant un mur. Noter $x$ la longueur d’un côté perpendiculaire au mur. Exprimer l’aire et trouver les dimensions optimales.}
\par\medskip
\textbf{Résolution.}\par
Les deux côtés perpendiculaires consomment $2x$ mètres ; le côté parallèle mesure $24-2x$. Les contraintes donnent $0<x<12$. L’aire vaut $A(x)=x(24-2x)=-2(x-6)^2+72$. Elle est au plus égale à $72$ m$^2$, valeur atteinte pour $x=6$. Le côté parallèle mesure alors $12$ m.
\end{minipage}\par
\par\noindent\begin{minipage}{\linewidth}\exercice{7}{Entraînement}{Un algorithme de résolution}
{\small\color{Muted}\textbf{Énoncé.} Décrire un algorithme qui reçoit $a,b,c$ avec $a\ne0$ et renvoie le nombre de racines réelles et leurs valeurs. L’appliquer à $a=1,b=0,c=-2$. Pourquoi traiter $\Delta=0$ séparément ?}
\par\medskip
\textbf{Résolution.}\par
Calculer $d=b^2-4ac$. Si $d<0$, annoncer aucune racine réelle. Si $d=0$, renvoyer la seule valeur $-b/(2a)$. Sinon, renvoyer $(-b-\sqrt d)/(2a)$ et $(-b+\sqrt d)/(2a)$. Pour $(1,0,-2)$, $d=8$ et les racines valent $\pm\sqrt2$. Le cas nul évite de présenter deux fois la même racine. Avec des données décimales, le test exact de nullité exige de prendre en compte les erreurs d’arrondi.
\end{minipage}\par
\par\noindent\begin{minipage}{\linewidth}\exercice{8}{Approfondissement}{Un paramètre}
{\small\color{Muted}\textbf{Énoncé.} Pour quels réels $m$ le trinôme $x^2-2mx+4$ est-il strictement positif pour tout réel $x$ ? Étudier les cas limites.}
\par\medskip
\textbf{Résolution.}\par
Le coefficient dominant est positif. La stricte positivité partout équivaut à l’absence de racine réelle, donc à $\Delta=4m^2-16<0$. On obtient $m^2<4$, soit $-2<m<2$. Pour $m=2$, le trinôme vaut $(x-2)^2$ et s’annule en $2$ ; pour $m=-2$, il vaut $(x+2)^2$ et s’annule en $-2$. Ces cas limites ne conviennent donc pas.
\end{minipage}\par
\end{document}
